%%%%%%%%%%%%%%%%%%%%%%%%%%%%
%%%%%     PREAMBLE     %%%%%
%%%%%%%%%%%%%%%%%%%%%%%%%%%%

\documentclass[11pt]{article}

%Define Margins using Geometry Package
\usepackage[top=1.0in,
bottom=1.0in,
right=1.0in,
left = 1.0in]{geometry}


%math
\usepackage{mathtools}


%Set up links and internal references
\usepackage[colorlinks=true,linkcolor=blue, urlcolor=blue]{hyperref} %hide the ugly red boxes
\usepackage[noabbrev]{cleveref} %don't abbreviate ``figure'', etc.
\usepackage{url} %for urls

%Remove natural indentation for itemize and enumerate environments
\usepackage{enumitem} %package to control itemize/enumerate behavior
\setlist[itemize]{leftmargin=*} %for itemize
\setlist[enumerate]{leftmargin=*} %for enumerate
%Swap out the itemize symbol for an N-dash rather than a bullet (does not require enumitem package)
%\def\labelitemi{--}

\usepackage{multicol} %Allow for two columns in the middle of a paper. 


\usepackage{titling}
\setlength{\droptitle}{-2.5cm}

%Set up fancy header and footer
\usepackage{fancyhdr}
\pagestyle{fancy}
%Fancy Header Content
\fancyfoot[L]{ME EN 497R}
\fancyfoot[C]{}
\fancyfoot[R]{Airframe Design Theory Exploration: Part II}
%Fancy Footer Content
% \fancyfoot[L,C]{} %make bottom left and center empty
% \fancyfoot[R]{\thepage}
\title{Airframe Design Theory Exploration: Part II}
\date{}


% Customize formatting for section headers
\renewcommand{\thesection}{Problem \arabic{section}}
\renewcommand{\thesubsection}{\arabic{section}.\alph{subsection}}

%%%%%%%%%%%%%%%%%%%%%%%%%%%%%%%%
%%%%%     END PREAMBLE     %%%%%
%%%%%%%%%%%%%%%%%%%%%%%%%%%%%%%%


%BEGIN Document Environment
\begin{document}
	
\maketitle
\vspace{-2cm}
	
%%%%%%%%%%%%%%%%%%%%%%%%%%%%%%%%%%%%%%%%%%%%%%%%%%%%%%%%%%%%%%%%%%%%%%
%                                                                    %
%                               Reading                              %
%                                                                    %
%%%%%%%%%%%%%%%%%%%%%%%%%%%%%%%%%%%%%%%%%%%%%%%%%%%%%%%%%%%%%%%%%%%%%%
\section*{Relevant Reading}
\label{sec:reading}

\begin{itemize}
	\item \href{http://flowlab.groups.et.byu.net/me415/flight.pdf}{Flight Vehicle Design Book}
	\begin{itemize}
		\item Chapter 4: Wing Design (7 pages)
	\end{itemize}
	\item \href{https://byu.box.com/shared/static/ywfayozbj3sr2ot6b32u8nqk5brqvurt.pdf}{Compuational Aeronautics}
	\begin{itemize}
		\item Chapter 4: Finite Wing (33 pages)
	\end{itemize}
\end{itemize}




%%%%%%%%%%%%%%%%%%%%%%%%%%%%%%%%%%%%%%%%%%%%%%%%%%%%%%%%%%%%%%%%%%%%%%
%                                                                    %
%                             Vocabulary                             %
%                                                                    %
%%%%%%%%%%%%%%%%%%%%%%%%%%%%%%%%%%%%%%%%%%%%%%%%%%%%%%%%%%%%%%%%%%%%%%
\section{Vocabulary}
\label{sec:vocab}

Explain the following terms; making sure to use sufficient detail, including any math or helpful figures.
In some cases, these terms are simple one sentence definitions, in others, you should include several paragraphs to explain them fully.

\paragraph{Non-dimensional Numbers}
\begin{itemize}
	\item Reynolds Number
	\item Mach Number
	\item Coefficients
	\begin{itemize}
		\item Lift Coefficient
		\item Drag Coefficient
		\item Moment Coefficient
	\end{itemize}
\end{itemize}

\paragraph{Airframe Performance}
\begin{itemize}
	\item Lift Distribution
	\item Stall Speed
\end{itemize}

%%%%%%%%%%%%%%%%%%%%%%%%%%%%%%%%%%%%%%%%%%%%%%%%%%%%%%%%%%%%%%%%%%%%%%
%                                                                    %
%                               Explore                              %
%                                                                    %
%%%%%%%%%%%%%%%%%%%%%%%%%%%%%%%%%%%%%%%%%%%%%%%%%%%%%%%%%%%%%%%%%%%%%%
\section{Theory}
\label{sec:theory}
In this problem, you get to dive into the theory the vortex lattice method. 

\begin{enumerate}[label=\roman*.]
	\item What is potential flow theory? What is the governing equation? What are the assumptions? 
	\item What is a vortex filament? What is a horseshoe vortex? 
	\item How do you use horseshoe vortices to model a lifting surface? What type of solver do you need to use? 
	\item Explain why an external drag model is required to capture stall.
\end{enumerate}


\section{Exploration}
\label{sec:exploration}

Complete the following exploration.

\subsection{Lift Distributions and Wing Efficiency}

\begin{enumerate}[label=\roman*.]
	\item Create a plot comparing the lift distributions of a constant chord, tapered chord, and elliptic wing with no twist and with the same wing area.  Also plot the ideal elliptic distribution for comparison.
	\item Discuss the differences you see in your plot.  Which planform design is the most efficient? Why?
\end{enumerate}




\end{document}